\chapter{Synthesis of the Predictive Controller}
\label{ch:deepc}

\section{Classical Model Predictive Control -- MPC}
\label{subs:mpc}

Classical MPC minimizes a receding-horizon cost function subject to system dynamics and actuator constraints \cite{nmpc}. For a linear time-invariant (LTI) system described by:
\begin{align}
    \boldsymbol{x}_{k+1} &= \boldsymbol{A}\boldsymbol{x}_k + \boldsymbol{B} \boldsymbol{u}_k 
    \label{eq:statempc}
    \\
    \boldsymbol{y}_k &= \boldsymbol{C} \boldsymbol{x}_k + \boldsymbol{D} \boldsymbol{u}_k
    \label{eq:outmpc}
\end{align}
where $\boldsymbol{x} \in \mathbb{R}^n$, $\boldsymbol{u} \in \mathbb{R}^m$, $\boldsymbol{y} \in \mathbb{R}^p$, the optimization problem over prediction horizon $N$ takes the form:
\begin{equation}
    \min_{\boldsymbol{u},\boldsymbol{x},\boldsymbol{y}} \sum_{k=0}^{N-1} \left( \left\| \boldsymbol{y}_k - \boldsymbol{r}_{k} \right\|_{\boldsymbol{Q}}^2 + \left\| \boldsymbol{u}_k \right\|_{\boldsymbol{R}}^2 \right)
    \label{eq:costmpc}
\end{equation}
subject to:
\begin{align}
    \boldsymbol{u}_k &\in \mathcal{U}, \quad \forall k \in \{0, \ldots, N-1\}
    \label{eq:umpc}
    \\
    \boldsymbol{y}_k &\in \mathcal{Y}, \quad \forall k \in \{0, \ldots, N-1\}
    \label{eq:ympc}
\end{align}

Despite its effectiveness, MPC requires a parametric model identified in a separate offline stage — a process that does not guarantee optimality from the control perspective \cite{markovsky2023data}. The DeePC algorithm eliminates this stage by operating directly on raw measurement data as a non-parametric behavioral representation of the system.

\section{Foundations of Behavioral Control Theory}

The foundation of the DeePC algorithm is the behavioral approach to systems theory, in which a dynamic system is defined as the set of all permissible input-output trajectories \cite{WILLEMS2005325}. Key concepts enabling such a representation are defined below \cite{quad, ethzurich, markovsky2023data, deepc}.

\textbf{Hankel Matrix} $\mathcal{H}$ -- for a discrete signal $\boldsymbol{w} = \begin{bmatrix} \boldsymbol{w}_1 & \boldsymbol{w}_2 & \dots & \boldsymbol{w}_T \end{bmatrix}^\mathrm{T}$ of length $T$, the Hankel matrix of depth $L = T_{ini} + N$ is defined as:
\begin{equation}
    \label{eq:hankel_matrix}
    \mathcal{H}_L(\boldsymbol{w}) = 
    \begin{bmatrix}
        \boldsymbol{w}_1 & \boldsymbol{w}_2 & \cdots & \boldsymbol{w}_{T-L+1}\\
        \boldsymbol{w}_2 & \boldsymbol{w}_3 & \cdots & \boldsymbol{w}_{T-L+2}\\
        \vdots & \vdots & \ddots & \vdots\\
        \boldsymbol{w}_L & \boldsymbol{w}_{L+1} & \cdots & \boldsymbol{w}_T
    \end{bmatrix}
\end{equation}

\textbf{Mosaic Hankel Matrix} -- for $K$ independent datasets $\boldsymbol{w}^{d_{1}}, \dots, \boldsymbol{w}^{d_{K}}$, the mosaic matrix is formed by horizontal concatenation:
\begin{equation}
    \label{eq:mosaic_hankel}
    \mathcal{H}_{mos}(\boldsymbol{w}^d) = 
    \begin{bmatrix}
        \mathcal{H}_L(\boldsymbol{w}^{d_{1}}) & \mathcal{H}_L(\boldsymbol{w}^{d_{2}}) & \cdots & \mathcal{H}_L(\boldsymbol{w}^{d_{K}})
    \end{bmatrix}
\end{equation}

\textbf{Persistency of Excitation (PE)} -- an input signal $\boldsymbol{u}$ is sufficiently exciting of order $L$ if $\mathcal{H}_L(\boldsymbol{u})$ has full row rank \cite{ethzurich}. For a system of order $n$, the input must be sufficiently exciting of order $T_{ini} + N + n$, requiring the total dataset length to satisfy:
\begin{equation}
    T \ge (m+1)(T_{ini} + N + n) - 1
\end{equation}

\textbf{Willems' Fundamental Lemma} -- for an LTI system with a sufficiently exciting input of order $T_{ini} + N + n$, the columns of $\mathcal{H}_L$ span the space of all possible input-output trajectories of length $L = T_{ini} + N$ \cite{deepc, WILLEMS2005325}. Consequently, any future system trajectory can be expressed as a linear combination of the columns of this matrix, without identifying the state matrices $\boldsymbol{A}, \boldsymbol{B}, \boldsymbol{C}, \boldsymbol{D}$.

\section{Data-Driven Control Algorithm -- DeePC}
\label{sec:deepc_algo}

The non-parametric model of the system is constructed from historical data $\boldsymbol{u}^d$, $\boldsymbol{y}^d$ of length $T$. The Hankel matrix is partitioned along the row dimension into past (index $p$, depth $T_{ini}$) and future (index $f$, depth $N$) blocks:
\begin{equation}
    \label{eq:data_matrix_deepc}
    \mathcal{H} = \begin{bmatrix}
        \mathcal{H}_{p} \\ \hline \mathcal{H}_{f}
    \end{bmatrix} = 
    \begin{bmatrix}
         \begin{pmatrix} \boldsymbol{U}_p \\ \boldsymbol{Y}_p \end{pmatrix} \\[1em]
         \hline
         \begin{pmatrix} \boldsymbol{U}_f \\ \boldsymbol{Y}_f \end{pmatrix}
    \end{bmatrix}
\end{equation}
as illustrated in Figure \ref{fig:deepc_matrix_structure}.

\begin{figure}[!htbp]
    \centering
    \includegraphics[width=0.6\linewidth]{rysunki/3.deepc/graf_deepc (3).drawio.pdf}
    \caption{Data matrix structure used in the DeePC algorithm. The gray area includes $T_{ini}$ rows responsible for establishing the initial state (past), while the white area contains $N$ rows defining possible future evolutions of the system. Each column represents one physically realizable system trajectory of length $L = T_{ini} + N$.}
    \label{fig:deepc_matrix_structure}
\end{figure}

\subsection{Trajectory Prediction and the Decision Variable $\boldsymbol{g}$}

A decision vector $\boldsymbol{g} \in \mathbb{R}^{T-L+1}$ is introduced, weighting the columns of the data matrix. At each control iteration, the sought trajectory must satisfy:
\begin{equation}
\label{eq:deePC_regression}
\begin{bmatrix}
    \boldsymbol{U}_p \\
    \boldsymbol{Y}_p \\
    \boldsymbol{U}_f \\
    \boldsymbol{Y}_f
    \end{bmatrix}
    \boldsymbol{g} =
    \begin{bmatrix}
    \boldsymbol{u}_{ini} \\
    \boldsymbol{y}_{ini} \\
    \boldsymbol{u} \\
    \boldsymbol{y}
    \end{bmatrix}
\end{equation}
where $\boldsymbol{u}_{ini}$, $\boldsymbol{y}_{ini}$ contain the last $T_{ini}$ measured samples (with $T_{ini} \ge n$ to uniquely define the initial state \cite{ethzurich}), and $\boldsymbol{u}$, $\boldsymbol{y}$ are the future control and output sequences over horizon $N$.

\subsection{Formulation of the Optimization Problem}

The DeePC control problem is formulated as a quadratic programming (QP) task:
\begin{equation}
    \label{eq:deepc_optimization}
    \begin{aligned}
    \min_{\boldsymbol{g}, \boldsymbol{u}, \boldsymbol{y}} \quad & \sum_{k=0}^{N-1} \left( \| \boldsymbol{y}_k - \boldsymbol{r}_{k} \|_{\boldsymbol{Q}}^2 + \| \boldsymbol{u}_k \|_{\boldsymbol{R}}^2 \right) \\
    \text{s.t.} \quad & \begin{bmatrix} \boldsymbol{U}_p \\ \boldsymbol{Y}_p \\ \boldsymbol{U}_f \\ \boldsymbol{Y}_f \end{bmatrix} \boldsymbol{g} = \begin{bmatrix} \boldsymbol{u}_{ini} \\ \boldsymbol{y}_{ini} \\ \boldsymbol{u} \\ \boldsymbol{y} \end{bmatrix} \\
    & \boldsymbol{u}_k \in \mathcal{U}, \quad \boldsymbol{y}_k \in \mathcal{Y}
    \end{aligned}
\end{equation}
For LTI systems without noise, this formulation is mathematically equivalent to classical MPC \cite{ethzurich}. The algorithm operates in receding horizon mode: at each step $t$, the optimization (\ref{eq:deepc_optimization}) is solved and only the first element $\boldsymbol{u}_{opt}(t) = \boldsymbol{u}^*_0 = \boldsymbol{U}_f\boldsymbol{g}^*$ is applied to the plant (Figure \ref{fig:algorytm_deepc}).

\begin{figure}[!htbp]
    \centering
    \includegraphics[width=1\linewidth]{rysunki/3.deepc/algorytm.drawio (2).pdf}
    \caption{Block diagram of the DeePC controller operation in a closed loop.}
    \label{fig:algorytm_deepc}
\end{figure}

\section{Robust Control for Noise and Nonlinearities -- Regularized DeePC}

To address measurement noise and plant nonlinearities, the \textbf{Regularized DeePC} variant \cite{deepc, ethzurich} introduces a slack variable $\boldsymbol{\sigma}$ and regularization terms into the cost function:
\begin{equation}
    \label{eq:regularized}
    \min_{\boldsymbol{g}, \boldsymbol{u}, \boldsymbol{y}, \boldsymbol{\sigma}} \sum_{k=0}^{N-1} \left( \| \boldsymbol{y}_k - \boldsymbol{r}_{k} \|_{\boldsymbol{Q}}^2 + \| \boldsymbol{u}_k \|_{\boldsymbol{R}}^2 \right) + \lambda_g \|\boldsymbol{g}\|_2^2 + \lambda_{\sigma} \|\boldsymbol{\sigma}\|^2_2
\end{equation}
where $\lambda_g$ prevents overfitting to measurement noise and improves numerical conditioning \cite{ethzurich, markovsky2023data}, and $\lambda_{\sigma}$ controls the degree of trust in current measurements. The slack variable modifies the initialization constraint:
\begin{equation}
\begin{bmatrix}
    \boldsymbol{U}_p \\
    \boldsymbol{Y}_p \\
    \boldsymbol{U}_f \\
    \boldsymbol{Y}_f
    \end{bmatrix}
    \boldsymbol{g} =
    \begin{bmatrix}
    \boldsymbol{u}_{ini} \\
    \boldsymbol{y}_{ini} + \boldsymbol{\sigma} \\
    \boldsymbol{u} \\
    \boldsymbol{y}
    \end{bmatrix}
    \label{eq:regularized_set}
\end{equation}
allowing the linear combination of historical data to reproduce the measurement $\boldsymbol{y}_{ini}$ within the accuracy of $\boldsymbol{\sigma}$. Quadratic ($l_2$) regularization was adopted, keeping the problem as a standard QP task while ensuring practical exponential stability in the presence of noise \cite{markovsky2023data}.

\section{Local Data Linearization -- Selective DeePC}
\label{sec:selective_deepc}

The \textbf{Selective DeePC} method \cite{choosewisely} addresses the fundamental limitation of Regularized DeePC when applied to highly nonlinear plants: operating on the full data library forces the algorithm to treat the system as a globally averaged LTI model, introducing significant modeling errors as the robot configuration changes. Its essence is the construction of a local non-parametric model at each control iteration by selecting only those columns of the data matrix that most closely resemble the current system state and the given control objective.

Denoting each column of the Hankel matrix as $\boldsymbol{h}_i$, the database vector and current state vector are defined as:
\begin{itemize}
    \item \textbf{Database vector} ($\boldsymbol{h}_i$) -- the $i$-th column of the Hankel matrix, consisting of historical inputs, outputs, and future outputs:
    \begin{equation}
        \boldsymbol{h}_i = 
        \begin{bmatrix}
        \boldsymbol{u}_{{ini}, i} \\
            \boldsymbol{y}_{{ini}, i} \\
            \boldsymbol{y}_{f, i}
        \end{bmatrix}
    \end{equation}
    
    \item \textbf{Current state vector} ($\boldsymbol{\tilde{h}}$) -- constructed at each control step from current measurements and the reference trajectory:
    \begin{equation}
        \boldsymbol{\tilde{h}} = 
        \begin{bmatrix}
            \boldsymbol{u}_{{ini},~{current}} \\
            \boldsymbol{y}_{{ini},~{current}} \\
            \boldsymbol{r}_{{future}}
        \end{bmatrix}
    \end{equation}
\end{itemize}

At each sampling step, the selection proceeds as follows:
\begin{enumerate}
    \item Calculation of the Euclidean distance $d_i$ between $\boldsymbol{\tilde{h}}$ and each column of the data library:
    \begin{equation}
        d_i = \| \boldsymbol{h}_i - \boldsymbol{\tilde{h}} \|, \quad \text{for } i = 1, \dots, M
    \end{equation}
    \item Sorting indices $i$ by increasing $d_i$.
    \item Selecting $N_{cols}$ nearest-neighbor columns best representing the current dynamics.
    \item Construction of a local Hankel matrix $\mathcal{\tilde H}$ from the selected columns.
\end{enumerate}

The matrix $\mathcal{\tilde H}$ is substituted into (\ref{eq:regularized}) in place of the full data matrix, reducing the dimension of $\boldsymbol{g}$ to $N_{cols}$ and improving control quality through local approximation of the nonlinear dynamics.